\chapter{The Dilemma of Agriculture}

The experience of the 1925 harvest (see p. 79 above) showed that the problem facing the planners of agricultural policy was not only to increase production, but to bring the product to the market; and it pointed ominously to the power of the well-to-do peasant and the kulak. The marketing crisis had, however, been surmounted, and optimism prevailed. The mood of hoping for the best was encouraged by the record harvest of 1926. When it had been gathered, many peasants had ample resources, and were sellers of grain. The market stringency of the previous year was not repeated, and prices were moderate. A feature of the grain collections was the increased participation of the agricultural trading cooperatives, which, though financed and controlled by the state, proved more popular and more efficient than the state purchasing agencies. The successful outcome of these operations helped to bring about the discomfiture of the opposition at the party conference in October 1926; the crisis predicted by the opposition had not occurred. It also encouraged the party leaders to step up programmes of industrialization, and to neglect the future consequences of this pressure for the peasant market. 
% 注：去除了 "jof", "/ample", "(party" 等前导的标点噪点。

In the external crises which overtook the USSR in 1927, and in the first flush of enthusiasm for planning, not much attention was paid to agriculture. The harvest, though lower than that of 1926, was satisfactory, and it was assumed that the grain collections would again proceed smoothly. This confidence was misplaced. The mood had changed since the previous year. The anxieties of the international situation, and talk of war and invasion, had spread to the countryside. After two good harvests the peasant was better off than at any time since the revolution. The well-to-do peasant had reserves both of grain and of money. The supply of industrial goods which he might want to buy was still meagre. The currency was once more being eroded by inflation; in a state of uncertainty and alarm, grain was the safest store of value. For those peasants who held stocks no incentive remained to put them on the market. The grain collections in the autumn of 1927, which should have been the best months, yielded less than half the quantities of 1926. But, if the well-to-do peasant had turned sour, provocation also came from the other side. The united opposition ever since its formation in the summer of 1926 had continued to denounce the policy of toleration for the kulak; and in October 1927 the party central committee, not to be outdone, called for ``a reinforced offensive against the kulak''. What happened in the autumn of 1927, though neither side perhaps at first realized its full import, was a declaration of war between the authorities and the well-to-do peasant who held the large available stocks of grain in the countryside. 
% 注：清除了极端的乱码 "iu ^ udu um __ uiii uii 0.11x1 K" ；去除了 "Iprevious", "'and talk", "iof value" 等乱码。

An atmosphere of false security prevailed at the party congress in December 1927. At the climax of the battle with the opposition, it would have been inopportune to admit that the country was in the throes of a grave crisis. Molotov observed regretfully that ``the advantage of the larger scale of production lies in practice on the side of the well-to-do peasant and the kulak''. But Stalin remarked mildly that ``the way out is to unite small and minute peasant holdings, gradually but surely, not by pressure, but by example and persuasion, into large holdings based on common, cooperative, collective cultivation of the land''; and the resolution of the congress, though it declared it to be ``the fundamental task of the party in the countryside'' to bring about ``the unification and conversion of independent peasant holdings into large collectives'', added that this could be done only with the consent of the ``labouring peasants''. Hard things were said at the congress about the kulaks. But the resolution went no further than to recommend higher and more progressive taxation of well-to-do peasants. No immediate emergency seemed to be pending. But no sooner was the congress over than the deadly nature of the threat to the food supplies of cities and factories was proclaimed in a series of decrees and emergency measures. Steps were taken---too late---to speed up the supply of textiles to country markets. Leading party members were despatched on tours of the principal grain-producing regions to supervise and enforce the collection of grain. Stalin made a three-weeks' tour of the main centres in western Siberia, where large stocks were believed to exist. 
% 注：去除了 "[jconsent", "II no", "11 taxation", "11 than", "'( of cities" 等多余字符。

What were called ``extraordinary measures'' were widely applied. An article of the criminal code was invoked which imposed the penalty of confiscation for concealment of grain. Propaganda and persuasion alternated with direct coercion. The situation was desperate. By hook or by crook, the holders of grain were induced to deliver it to the collecting organs; it was admitted that the recalcitrant holders of grain were not all kulaks, and that many so-called ``middle peasants'' were also compelled to disgorge their reserves. Not much distinguished these procedures from the wholesale requisitions of the days of war communism. Between January and March 1928 very large quantities of grain were obtained, and in March Rykov announced that the grain crisis had been ``taken off the agenda''. The first battle of the grain had been won by the government, but in conditions which promised that the war would continue, and would be fought out with extreme bitterness. On one side, well-to-do peasants had been harshly, often brutally, treated. On the other, bread queues appeared in the cities; and scarce foreign currency, sorely needed to finance industrialization, had to be spent on imports of grain to make up the deficiency. No easy or acceptable answer could be found to the question who should suffer from the grain shortage. 
% 注：去除了 "fill What", "rmapplied", "^|fPropaganda", "^flThe", "A harshly", "-7I of grain" 等边缘噪点，并将 "ot all kulaks" 修正为 "not all kulaks"，修正 "Were also" 为 "were also"。

The harshness of the ``extraordinary measures'' shook and divided the party. Many workers still had close ties with the countryside, and were well aware of what had been done. Disaffection was said to have spread to the ranks of the Red Army, composed predominantly of peasants. Rykov was apparently the first of the party leaders to express disquiet, and was soon joined by Bukharin, throughout the NEP period the leading advocate of the conciliation of the peasant, and Tomsky, who was now seriously disturbed by the pressures imposed by industrialization on the workers and on the trade unions. At a crucial session of the party central committee in July 1928 the line was drawn between those who wished to relax the pressure on the peasant, even at the cost of slowing down the pace of industrialization, and those who gave unconditional priority to industrialization, however severe the measures of coercion inflicted on the peasant. Bukharin spoke of ``a wave of dissatisfaction'' and ``outbreaks'' in the countryside, and of ``a return to war communism''. Stalin professed to temporize, admitted that excesses had occurred, and expressed the belief that these would be avoided in the forthcoming grain collections. Some increase in agricultural prices was conceded. The resolution condemned ``violations of revolutionary legality'' and ``the frequent application of methods of requisition''. It was a hollow compromise. No doubt remained that the party machine, with Stalin, Molotov and Kuibyshev in command, was now firmly geared to all-out industrialization, and that whatever measures were required to secure supplies of food for the cities and the workers would be taken. 
% 注：去除了 "I ", "i[ done", "^/(the" 和 "[slowing" 的乱码。

The experience of the grain collections in the autumn of 1928 repeated that of the previous year on a larger scale. The total figure for the harvest had been maintained. But the essential crops for human consumption, wheat and rye, had fallen off. The authorities were more aware than before of the critical situation, and were more determined and ruthless. Organization had improved; a new central organ, Soyuzkhleb, was set up to control the collections. The reserve stocks in the hands of the peasants had been depleted by the raids of the spring of 1928. The peasants were better prepared for a fresh onslaught, and more adept at concealing what they held. More important still, stringency in the cities revived and expanded the black market. Private traders travelled far and wide in rural areas, offering prices for grain much in excess of the modestly increased official prices. The battle on both sides was fiercely engaged. Legal pretexts were once more invoked to justify confiscation. Reprisals were frequent for real or imaginary offences. In the Urals and in Siberia a system was devised by which the village Soviet or the village assembly was induced to agree on a quota for the village, which was then imposed on the well-to-do peasant under threat of severe penalties. This system, nicknamed the ``Ural-Siberian method'' of grain collection, was subsequently extended to other regions, and was a powerful instrument of victimization. Disorders and mass demonstrations against the grain collections were reported. Greater latitude was apparently allowed to the local authorities to make ``decentralized'' collections of grain for delivery to towns within the region, so that comparison with the previous year's figures may be misleading. But only 8.3 million tons of grain were collected by the central collecting organs in 1928-1929 against 10.3 millions in 1927-1928; and this included only 5.3 million tons of wheat and rye against 8.2 millions in 1927-1928. What was certain was that nobody delivered grain to the official agencies except under some degree of coercion or fear. 
% 注：去除了段首错落的 "[The", "[the", "[had", "jWhat"。将表示小数的 "8-3"、"10-3" 等格式修复为了标准的英文浮点数 "8.3"、"10.3"。

The situation quickly became desperate. Bread was the staple diet of peasant and industrial worker alike; Moscow, with a population less than two-fifths that of Berlin, consumed more bread. With the progress of industrialization the urban population was rapidly increasing. In the winter of 1927-1928 bread queues were familiar in the cities; and butter, cheese and milk were rarely to be had. State stocks of grain were exhausted. The deficiency was made good by importing grain, by mixing barley and maize with wheat and rye in bread-making, by coarse grinding, by bread rationing in the large cities, and by recourse to the black market. The restraints imposed on the private trader in the towns were inoperative in the countryside. Even attempts to discriminate against him in the provision of such facilities as storage and transport proved largely ineffective. Of the marketed grain from the 1928 harvest 23 per cent was taken by private traders. In July 1928 the party central committee, faced with this competition and with the demands of the new ``Right'' opposition in the party, had decided to concede an increase in the official prices for grain. The prices paid for grain by the official collecting agencies in 1928-1929 were about 20 per cent higher than in the previous year. But prices on the black market soared faster. It was calculated that the prices of agricultural products in private trade, which in 1927-1928 exceeded the official prices by about 40 per cent, were almost double the official prices in the following year. This may well have been an under-estimate. The scarcity was now chronic. Bread cards were introduced in Leningrad in November 1928 and in Moscow in March 1929; these were honoured up to the extent of available supplies. No cards were issued to non-workers, who were left to the mercy of the private market. Belts had to be tightened all round. Nor was there any security for the future. The belief, on which NEP was founded, that the cities could be fed through a combined system of voluntary deliveries to the state and free sales on the market had broken down. 
% 注：修复了极其严重的多栏交错断句 ("...faced with this competi- »in the party, had decided to concede an increase in the tion and with the demands of the new “Right” opposition official prices for grain.") 并理顺了原本连续的英语逻辑。

The dilemma brought into the open a fundamental issue hitherto hidden beneath the surface of day-to-day politics. The maintenance of small individual peasant holdings, often divided into widely separated strips, and practising the ancient three-field rotation, was plainly incompatible with an efficient agriculture. Even in a primitive peasant economy some measure of cooperation between cultivators of the soil was demanded by common sense as well as by socialist doctrine. In pre-revolutionary Russia, this was provided over the greater part of the country at two levels. The smallest unit of production was not the individual, but the household or dvor. This varied in size, having in the past often been an ``extended'' family comprising more than one generation, sometimes supplemented by adoption. In most parts of Russia a group of dvors formed a land-holding community (mir or obshchina) which regulated matters of common concern such as pasture, rotation of crops, ditching, fencing and road-building, and in most cases periodically redistributed the land to take account of the changing composition of the dvors. 
% 注：去除了 "r", "j", "1", "•r" 等侧面乱码。

After the revolution the dvors increased in number, but declined in importance, owing to the frequent splitting of family units. This was said to be due to the greater insistence on independence by the younger generation, and especially by the women, who on marriage demanded the establishment of a separate household. But the traditional authority of the mir, strengthened by the disappearance of the landlord and the weakening of the dvor, increased; it was often a successful rival of the newly established village Soviet. The official attitude to the mir was ambivalent and inconsistent. On the one hand, the mir generally resisted attempts to change such ancient agricultural practices as the three-field rotation or the strip system of land-holding. Sometimes the enterprising peasant, or the kulak, was said to dominate the other peasants in the mir for his own advantage. Sometimes he left the mir, taking his share of land with him, and set up an independent unit of cultivation. In all these ways, the mir perpetuated traditions of the past which the revolution aimed to eradicate. On the other hand, the mir was by far the most effective existing agency of collective action in a peasant community. Narodniks like Herzen regarded the mir as a stepping-stone to socialism; and Marx could be quoted (though his very tentative comment related to a now remote period) as contemplating the possibility that, in the event of ``a workers' revolution in the west'', the Russian mir might serve as ``a starting-point for communist development''. During the nineteen-twenties, the status of the mir was much debated in Moscow. But no attempt was made to interfere with it on the spot; and it survived more or less untouched till the moment of collectivization. The power of party and government in the countryside was still extremely weak. 
% 注：此处解决了整章最严重的跨栏交织断句。将 "taking his share..." 与前文正常衔接；将 "On the other hand, the mir was by far the most effective existing agency..." 和 "Narodniks like Herzen regarded..." 剥离并恢复正常顺序。去除了大量伴生的乱码 ("III taking", "H/l'", "* 1/^^", "f Narodniks", "1/ existing", "-^/(/countryside")。

Forms of organization designed to promote collective cultivation were mainly the creation of the Soviet period. Agricultural cooperatives had flourished before the revolution, and continued to exist, but were concerned with collective marketing, and the provision of credit and the purchase of machines, rather than with collective production. Collective farms (Kolkhozy) and Soviet farms (Sovkhozy), which dated from the days of war communism (see pp. 22-23 above), wilted under NEP. They were said to occupy less than 2 per cent of all land in the USSR, and for some years received little support from the authorities. Many Kolkhozy survived as loose cooperatives, with the collective principle enjoying a much diminished respect. Many Sovkhozy became a by-word for inefficiency. The grain crises of the middle nineteen-twenties led to a renewal of interest in both institutions. The large collective unit was more likely to provide surpluses for the market than the individual peasant working primarily for the needs of himself and his family, and was more amenable to the inducements and pressures of the grain collections. The multiplication of small peasant holdings after the revolution had aggravated the stringency. While Bukharin and his disciples still pinned their faith on the agricultural marketing cooperatives, and hoped that these would lead gradually to an extension of collective cultivation, official policy began to swing towards a revival of the Kolkhozy. A central organization, Kolkhoztsentr, was set up in 1926. A movement began to form new Kolkhozy, which multiplied rapidly after the middle of 1927. These were smaller than the original Kolkhozy of the period of war communism; and the practice of collective working among their members did not extend very far. But they were a significant attempt to overcome the traditional conservatism of the mass of peasants, as well as the interested opposition of the prosperous kulak. 
% 注：去除了 "^ support", "//loose", "-;^//much", "if /collective", "■jIaJ the needs"。修复了分页断词 "organiza- don" 为 "organization"。

The Sovkhozy lagged behind the Kolkhozy; their revival did not begin before 1927, and was connected with the process sometimes designated by the catchword ``the industrialization of agriculture''. The substitution of even the simplest machines and implements for the primitive tools of the traditional Russian peasant (typified by the wooden plough, which he could make for himself) had long been recognized as a crying need, and attempts had been made to meet it through agricultural credit cooperatives and a state-financed agricultural bank. More ambitiously, Lenin had proclaimed that ``100,000 first-class tractors'' were required to convert the peasant to communism. In the early nineteen-twenties the Putilov works in Leningrad built a few tractors on an American model; and from 1923 onwards some hundreds of tractors were imported from the United States. In 1925 a plan was first mooted to build a large tractor factory at Stalingrad; it was three years before it was finally approved, and a start made on the work. Official propaganda underlined the intended role of the Sovkhozy as model farms, not only setting the example of modern methods of cultivation for the surrounding peasant holdings, but providing them with tractors and other agricultural machines. These ideals were far from being fulfilled. Yet it was becoming clear that, if agriculture was to be made more efficient through the use of tractors and complex machines, the work must be conducted in larger units than the individual peasant holding. In spite of party exhortations, little progress had yet been made either with mechanization or with collectivization. This was to be the task of the next period.
% 注：去除了 "/sometimes", "/of agriculture", "/need", "' first-class", "in- j tended" 等乱码，并修复了断字。